% Generated draft: requires mathematical and source review.
\documentclass[11pt]{article}
\usepackage[margin=25mm]{geometry}
\usepackage{fontspec}
\setmainfont{Latin Modern Roman}
\usepackage{amsmath,amssymb,mathtools,unicode-math}
\setmathfont{Latin Modern Math}
\usepackage{enumitem,needspace,xurl,hyperref}
\setlength{\parindent}{0pt}
\setlength{\parskip}{5pt}
\setlength{\emergencystretch}{3em}
\newcommand{\N}{\mathbb{N}}
\newcommand{\Z}{\mathbb{Z}}
\newcommand{\Q}{\mathbb{Q}}
\newcommand{\R}{\mathbb{R}}
\newcommand{\C}{\mathbb{C}}
\newcommand{\F}{\mathbb{F}}
\newcommand{\A}{\mathbb{A}}
\newcommand{\PP}{\mathbb{P}}
\newcommand{\E}{\mathbb{E}}
\newcommand{\eps}{\varepsilon}
\newcommand{\dd}{\,\mathrm{d}}
\begin{document}
\section*{UnsolvedMath ID: 30004312}
\textbf{Source problem code:} catalogue entry
\subsection{Definitions and mathematical statement}
A \emph{skew left brace} is a set \(A\) with group operations \(+\) and \(\circ\), having the same identity \(0\), such that
\[
a\circ(b+c)=(a\circ b)-a+(a\circ c)
\qquad(a,b,c\in A).
\]
The additive group need not be abelian. Define
\[
\lambda_a(b)=-a+(a\circ b),\qquad a*b=\lambda_a(b)-b.
\]
A subbrace is a subset that is a subgroup for both operations. A brace is \emph{trivial} if \(a\circ b=a+b\) for every pair of its elements; this does not require its additive group to be abelian. Using the right-series convention
\[
A^{(1)}=A,\qquad
A^{(n+1)}=\bigl\langle x*y:x\in A^{(n)},\ y\in A\bigr\rangle_{+}
\quad(n\geq1),
\]
where \(\langle\cdot\rangle_{+}\) denotes the subgroup generated in \((A,+)\), call \(A\) \emph{right-nilpotent} if \(A^{(m)}=\{0\}\) for some integer \(m\geq1\).

\emph{First question, with an explicit conditional interpretation of ``sum.''} If the source means an additive set factorization by subbraces, its question becomes: for every skew left brace \(A\) and subbraces \(B,C\subseteq A\), does
\[
A=B+C:=\{b+c:b\in B,\ c\in C\},\qquad
b\circ b'=b+b'\quad(b,b'\in B),\qquad
C^{(m)}=\{0\}\text{ for some }m\geq1
\]
imply
\[
A^{(n)}=\{0\}\text{ for some }n\geq1?
\]
Here the series for \(C\) uses its inherited brace operations. This is a conditional formulation, not a determination of the source's intended meaning: the supplied record does not define ``sum,'' require the summands to be ideals, specify a multiplicative factorization, or impose finiteness in this first question. These missing qualifications must be resolved before treating this displayed implication as the definitive source statement.

\emph{Further adaptation questions.} Determine whether there are appropriate skew-left-brace versions of the Kegel--Wielandt theorem and of the result attributed in the record to Huppert concerning finite sums of finite cyclic groups and supersolvability. The classical Kegel--Wielandt theorem concerns a finite group that is a set product of two nilpotent subgroups and concludes that the group is solvable. The supplied record does not specify the corresponding brace hypotheses or the intended brace notion of solvability. For the Huppert question, it also does not specify the number of cyclic factors, the meaning of their sum, whether cyclicity concerns the additive or multiplicative groups, or a brace notion of supersolvability. Consequently, these requests do not yet determine unique quantified conjectures. In particular, the phrase ``finite sums of finite cyclic groups'' cannot safely be interpreted as an unrestricted finite product of cyclic subgroups: every finite group is generated by finitely many cyclic subgroups, and such a supersolvability assertion would require additional restrictions.

Finally, determine what consequences the intended brace notions have for associated Yang--Baxter solutions. For reference, a bijective set-theoretic solution is a set \(X\) and a bijection \(r:X\times X\to X\times X\) satisfying
\[
(r\times\operatorname{id}_X)(\operatorname{id}_X\times r)(r\times\operatorname{id}_X)
=
(\operatorname{id}_X\times r)(r\times\operatorname{id}_X)(\operatorname{id}_X\times r)
\quad\text{on }X^3.
\]
The record does not identify a particular solution property, additional solution hypotheses, or a precise implication to establish.

The supplied dated assessment reports partial results for strong trivial summands and right-nilpotency bounds, and examples obstructing a naive Kegel--Wielandt analogue. It describes the full sum and adaptation questions as remaining open. These are supplied literature qualifications, not independently verified mathematical conclusions; neither the extra hypotheses in those partial results nor the obstructed naive formulation are specified here.
\subsection{Short English statement}
The problem asks whether a sum of a trivial skew left brace and a right-nilpotent skew left brace must be right-nilpotent, and seeks brace analogues of Kegel--Wielandt and Huppert results together with consequences for Yang--Baxter solutions. The meaning of sum and several hypotheses and target notions remain unspecified. The supplied classification is partially solved; its dated literature assessment requires independent review.
\subsection{Sources}
\begin{itemize}
\item Supplied local UnsolvedMath registry record, permanent ID \(30004312\), title ``Kegel--Wielandt Theory for Skew Braces,'' source problem code \texttt{catalogue entry}, in the Oberwolfach Reports Open Problems collection. The record gives proposed year 2019 and names no individual proposer.
\item Supplied report citation: \emph{Mini-Workshop: Algebraic Tools for Solving the Yang--Baxter Equation} (2020), Oberwolfach Reports, pp.\ 3207--3242, DOI 10.4171/owr/2019/51. \url{https://doi.org/10.4171/owr/2019/51}. This is supplied report-level provenance; no problem-specific page or passage is provided.
\item Supplied literature source: \emph{Factorizations of skew braces}. \url{https://notablesdelaciencia.conicet.gov.ar/bitstream/handle/11336/136170/CONICET_Digital_Nro.11b9658a-b246-481f-ae3c-c31eaf455ce5_X.pdf?isAllowed=y\&sequence=5}.
\item Supplied literature source: \emph{A generalization of Ito's theorem to skew braces}. \url{https://www.sciencedirect.com/science/article/pii/S0021869323006233}.
\item Supplied literature source: \emph{Oberwolfach Reports, skew braces problems}. \url{https://ems.press/content/serial-article-files/46831}.
\item Literature assessment date: \texttt{2026-08-22}. This is supplied dated curation, not independently verified. The assessment labels the problem partially solved, reports factorization results for strong trivial summands and right-nilpotency bounds, reports examples against a naive Kegel--Wielandt analogue, and describes the full sum and adaptation questions as open. The background explicitly requests independent verification of time-sensitive status and cited claims.
\item The registry labels its statement \texttt{corrected\_verified} and reports a manual semantic, typographic, and LaTeX review. These are supplied registry assertions and do not constitute independent verification of the present mathematical formulation.
\end{itemize}
\end{document}
